%
En beholder inneholder en blanding av heliumgass og oksygengass med temperatur \SI{300}{\kelvin}. Massen til et heliumatom er \SI{6.64e-27}{\kilo\gram}, og massen til et oksygenmolekyl er \SI{5.31e-26}{\kilo\gram}.
\begin{enumerate}[a)]
    \item Hvilken av de to har størst gjennomsnittlig kinetisk energi?
    \item Bruk $\bar{E}_\text{K} = \tfrac{1}{2}mv^2$ til å finne en typisk fart for heliumatomene og for oksygenmolekylene.
    \item Hvor mange ganger raskere beveger heliumatomene seg?
\end{enumerate}
%
